% Shared preamble for the pyMack LST documentation set.
% Used by: mean_boundary_layer.tex, disturbance_equations.tex,
%          numerical_methods.tex, validation.tex
% Each document begins with:  \documentclass[11pt]{article}% Shared preamble for the pyMack LST documentation set.
% Used by: mean_boundary_layer.tex, disturbance_equations.tex,
%          numerical_methods.tex, validation.tex
% Each document begins with:  \documentclass[11pt]{article}% Shared preamble for the pyMack LST documentation set.
% Used by: mean_boundary_layer.tex, disturbance_equations.tex,
%          numerical_methods.tex, validation.tex
% Each document begins with:  \documentclass[11pt]{article}% Shared preamble for the pyMack LST documentation set.
% Used by: mean_boundary_layer.tex, disturbance_equations.tex,
%          numerical_methods.tex, validation.tex
% Each document begins with:  \documentclass[11pt]{article}\input{lst_preamble}
\usepackage[margin=1in]{geometry}
\usepackage{amsmath,amssymb,mathtools}
\usepackage{booktabs}
\usepackage{enumitem}
\usepackage{array}
\usepackage{graphicx}
\usepackage{caption}
\usepackage{xcolor}
\usepackage[hidelinks]{hyperref}

\graphicspath{{figures/}}
\captionsetup{font=small,labelfont=bf}

\definecolor{kicol}{RGB}{31,143,107}
\definecolor{boxbg}{RGB}{238,246,242}
\definecolor{rmkbg}{RGB}{244,244,244}

\newcommand{\plain}[1]{\par\smallskip\noindent{\color{kicol}\textbf{In plain words.}}\ \textit{#1}\par\smallskip}
\newcommand{\problembox}[1]{%
  \par\smallskip\noindent\fcolorbox{kicol}{boxbg}{%
  \parbox{0.965\linewidth}{\smallskip\textbf{\color{kicol}What problem is solved here.}\ \ #1\smallskip}}\par\smallskip}
\newcommand{\remark}[2]{%
  \par\smallskip\noindent\colorbox{rmkbg}{%
  \parbox{0.965\linewidth}{\smallskip\footnotesize\textbf{#1}\ \ #2\smallskip}}\par\smallskip}
\newcommand{\bridge}[1]{\par\noindent\textit{#1}\par\medskip}

\newcommand{\D}{\mathrm{D}}
\newcommand{\imag}{\mathrm{i}}
\newcommand{\Real}{\operatorname{Re}}
\newcommand{\Imag}{\operatorname{Im}}
\newcommand{\Mae}{M_e}
\newcommand{\Ree}{\mathit{Re}}
\renewcommand{\bar}[1]{\overline{#1}}
\newcommand{\hatu}{\hat{u}}
\newcommand{\hatv}{\hat{v}}
\newcommand{\hatT}{\hat{T}}
\newcommand{\hatp}{\hat{p}}
\newcommand{\hatq}{\hat{q}}
\newcommand{\nuR}{\nu_{\!R}}

\usepackage[margin=1in]{geometry}
\usepackage{amsmath,amssymb,mathtools}
\usepackage{booktabs}
\usepackage{enumitem}
\usepackage{array}
\usepackage{graphicx}
\usepackage{caption}
\usepackage{xcolor}
\usepackage[hidelinks]{hyperref}

\graphicspath{{figures/}}
\captionsetup{font=small,labelfont=bf}

\definecolor{kicol}{RGB}{31,143,107}
\definecolor{boxbg}{RGB}{238,246,242}
\definecolor{rmkbg}{RGB}{244,244,244}

\newcommand{\plain}[1]{\par\smallskip\noindent{\color{kicol}\textbf{In plain words.}}\ \textit{#1}\par\smallskip}
\newcommand{\problembox}[1]{%
  \par\smallskip\noindent\fcolorbox{kicol}{boxbg}{%
  \parbox{0.965\linewidth}{\smallskip\textbf{\color{kicol}What problem is solved here.}\ \ #1\smallskip}}\par\smallskip}
\newcommand{\remark}[2]{%
  \par\smallskip\noindent\colorbox{rmkbg}{%
  \parbox{0.965\linewidth}{\smallskip\footnotesize\textbf{#1}\ \ #2\smallskip}}\par\smallskip}
\newcommand{\bridge}[1]{\par\noindent\textit{#1}\par\medskip}

\newcommand{\D}{\mathrm{D}}
\newcommand{\imag}{\mathrm{i}}
\newcommand{\Real}{\operatorname{Re}}
\newcommand{\Imag}{\operatorname{Im}}
\newcommand{\Mae}{M_e}
\newcommand{\Ree}{\mathit{Re}}
\renewcommand{\bar}[1]{\overline{#1}}
\newcommand{\hatu}{\hat{u}}
\newcommand{\hatv}{\hat{v}}
\newcommand{\hatT}{\hat{T}}
\newcommand{\hatp}{\hat{p}}
\newcommand{\hatq}{\hat{q}}
\newcommand{\nuR}{\nu_{\!R}}

\usepackage[margin=1in]{geometry}
\usepackage{amsmath,amssymb,mathtools}
\usepackage{booktabs}
\usepackage{enumitem}
\usepackage{array}
\usepackage{graphicx}
\usepackage{caption}
\usepackage{xcolor}
\usepackage[hidelinks]{hyperref}

\graphicspath{{figures/}}
\captionsetup{font=small,labelfont=bf}

\definecolor{kicol}{RGB}{31,143,107}
\definecolor{boxbg}{RGB}{238,246,242}
\definecolor{rmkbg}{RGB}{244,244,244}

\newcommand{\plain}[1]{\par\smallskip\noindent{\color{kicol}\textbf{In plain words.}}\ \textit{#1}\par\smallskip}
\newcommand{\problembox}[1]{%
  \par\smallskip\noindent\fcolorbox{kicol}{boxbg}{%
  \parbox{0.965\linewidth}{\smallskip\textbf{\color{kicol}What problem is solved here.}\ \ #1\smallskip}}\par\smallskip}
\newcommand{\remark}[2]{%
  \par\smallskip\noindent\colorbox{rmkbg}{%
  \parbox{0.965\linewidth}{\smallskip\footnotesize\textbf{#1}\ \ #2\smallskip}}\par\smallskip}
\newcommand{\bridge}[1]{\par\noindent\textit{#1}\par\medskip}

\newcommand{\D}{\mathrm{D}}
\newcommand{\imag}{\mathrm{i}}
\newcommand{\Real}{\operatorname{Re}}
\newcommand{\Imag}{\operatorname{Im}}
\newcommand{\Mae}{M_e}
\newcommand{\Ree}{\mathit{Re}}
\renewcommand{\bar}[1]{\overline{#1}}
\newcommand{\hatu}{\hat{u}}
\newcommand{\hatv}{\hat{v}}
\newcommand{\hatT}{\hat{T}}
\newcommand{\hatp}{\hat{p}}
\newcommand{\hatq}{\hat{q}}
\newcommand{\nuR}{\nu_{\!R}}

\usepackage[margin=1in]{geometry}
\usepackage{amsmath,amssymb,mathtools}
\usepackage{booktabs}
\usepackage{enumitem}
\usepackage{array}
\usepackage{graphicx}
\usepackage{caption}
\usepackage{xcolor}
\usepackage[hidelinks]{hyperref}

\graphicspath{{figures/}}
\captionsetup{font=small,labelfont=bf}

\definecolor{kicol}{RGB}{31,143,107}
\definecolor{boxbg}{RGB}{238,246,242}
\definecolor{rmkbg}{RGB}{244,244,244}

\newcommand{\plain}[1]{\par\smallskip\noindent{\color{kicol}\textbf{In plain words.}}\ \textit{#1}\par\smallskip}
\newcommand{\problembox}[1]{%
  \par\smallskip\noindent\fcolorbox{kicol}{boxbg}{%
  \parbox{0.965\linewidth}{\smallskip\textbf{\color{kicol}What problem is solved here.}\ \ #1\smallskip}}\par\smallskip}
\newcommand{\remark}[2]{%
  \par\smallskip\noindent\colorbox{rmkbg}{%
  \parbox{0.965\linewidth}{\smallskip\footnotesize\textbf{#1}\ \ #2\smallskip}}\par\smallskip}
\newcommand{\bridge}[1]{\par\noindent\textit{#1}\par\medskip}

\newcommand{\D}{\mathrm{D}}
\newcommand{\imag}{\mathrm{i}}
\newcommand{\Real}{\operatorname{Re}}
\newcommand{\Imag}{\operatorname{Im}}
\newcommand{\Mae}{M_e}
\newcommand{\Ree}{\mathit{Re}}
\renewcommand{\bar}[1]{\overline{#1}}
\newcommand{\hatu}{\hat{u}}
\newcommand{\hatv}{\hat{v}}
\newcommand{\hatT}{\hat{T}}
\newcommand{\hatp}{\hat{p}}
\newcommand{\hatq}{\hat{q}}
\newcommand{\nuR}{\nu_{\!R}}
